\documentclass[11pt]{article}

\usepackage[margin=1in]{geometry}
\usepackage{graphicx}
\usepackage{booktabs}
\usepackage{float}
\usepackage{hyperref}
\usepackage{xcolor}

\hypersetup{
    colorlinks=true,
    linkcolor=blue,
    urlcolor=blue,
    citecolor=blue
}

% ---------- Replace these values ----------
\newcommand{\studentname}{Annsh Navle, Rohan Malpani, Valmik Vyas}
\newcommand{\repolink}{https://github.com/RohMalpani/CS422Assignments}

% Use this command for links to specific files or line ranges in the repository.
% Example: \codelink{Ping collection code}{src/collect.py\#L20-L75}
\newcommand{\codelink}[2]{\href{\repolink/blob/main/#2}{#1}}

% Displays an image if it exists; otherwise displays a labeled placeholder.
\newcommand{\plotplaceholder}[2]{%
    \IfFileExists{#1}{%
        \includegraphics[width=0.88\textwidth]{#1}%
    }{%
        \fbox{\parbox[c][2.7in][c]{0.82\textwidth}{\centering
        Plot placeholder\\[0.5em]
        Save the finished plot as \texttt{\detokenize{#1}}}}%
    }%
}

\title{Assignment 1: Network Latencies, Ping, and Traceroute}
\author{\studentname}
\date{CS 422 -- Fall 2026}

\begin{document}

\maketitle

\section*{Repository}

The complete source code, input data, and instructions needed to reproduce the
experiments and plots are available at:
\begin{center}
    \url{\repolink}
\end{center}

% Briefly state the command used to reproduce all results.
\noindent Reproduction command:
\begin{verbatim}
./run_all.sh
\end{verbatim}

\section{Ping Test and Round-Trip Time}

\subsection{Method}

We tested every destination listed in the provided IP-address file, along with
our own public IP address.  The script sent ten ICMP echo requests to each
destination and recorded the minimum, average, and maximum RTT, using
\texttt{NA} for non-responsive destinations.  It obtained latitude and
longitude from \url{ip-api.com} and calculated great-circle distance with the
Haversine formula using an Earth radius of 6,371 km.

The code used to collect and process the ping measurements is available here:
\codelink{ping collection}{assignment1/part1/ping.sh\#L19-L20} and
\codelink{RTT calculation}{assignment1/part1/ping.sh\#L22-L30}.

\subsection{Results}

The first five rows are shown below (showing 5/37 results only).  The complete
data are available in \texttt{part1/results.csv}.

\begin{table}[H]
    \centering
    \caption{First five ping results.}
    \label{tab:ping-results}
    \begin{tabular}{lrrrr}
        \toprule
        Destination IP & Distance (km) & Min RTT (ms) & Avg RTT (ms) & Max RTT (ms) \\
        \midrule
        52.119.103.17  & 0      & 4.209   & 12.189  & 56.938  \\
        160.242.19.254 & 11,498 & 298.388 & 319.136 & 367.114 \\
        41.110.39.130  & 7,463  & 164.185 & 182.277 & 246.594 \\
        213.158.175.240 & 9,918 & 168.356 & 222.082 & 356.200 \\
        102.214.66.39  & 9,327  & 211.874 & 248.932 & 322.786 \\
        \bottomrule
    \end{tabular}
\end{table}

\subsection{Distance versus RTT}

\begin{figure}[H]
    \centering
    \plotplaceholder{figures/distance_vs_rtt.pdf}{Distance versus RTT}
    \caption{Geographic distance from West Lafayette, Indiana, versus
    round-trip time. Each point represents one destination IP address.}
    \label{fig:distance-rtt}
\end{figure}

The plotting code is available here:
\codelink{distance-versus-RTT plot}{assignment1/part1/scatter.py\#L1-L74}.

\subsection{Discussion}

The measurements show a moderate positive relationship between geographic
distance and RTT.  Among the 35 destinations with both valid geolocation and
ping data, the Pearson correlation between distance and average RTT is
approximately 0.57.  The correlations for minimum and maximum RTT are similar,
at approximately 0.58 and 0.57, respectively.  Thus, more distant destinations
generally have greater RTTs, but distance alone does not predict latency
precisely.  Geographic distance is the great-circle distance between endpoints,
whereas packets follow routes selected by network providers.  A route may be
longer than the direct geographic path because of peering locations, submarine
cable landing points, or indirect routing.  Every router along that route also
adds processing and transmission delay.

The minimum RTT is the best estimate in this experiment of the path's baseline
latency.  It reflects propagation delay plus the smallest observed amount of
router processing, transmission, and queueing delay.  It does not measure
distance directly.  Destinations at the same reported location can still have
different minimum RTTs because their providers may use different routes.  For
example, the four Tokyo-area servers are all about 10,304 km away, but their
minimum RTTs range from approximately 135 ms to 186 ms.  This variation is
consistent with differences in routing and network infrastructure rather than
geographic distance.

The maximum RTT describes the slowest of the ten ping replies and is more
sensitive to the network state during the measurement.  Temporary congestion
can increase queueing delay, and routers may deprioritize ICMP echo traffic.
For example, \texttt{speedtestfl.telecom.mu} had a minimum RTT of approximately
311 ms but a maximum of approximately 662 ms, a spread of about 351 ms.  The
Singapore destination \texttt{96.45.38.22} ranged from approximately 203 ms to
488 ms.  These large spreads cannot be explained by a change in physical
distance during the test.  They indicate transient queueing, congestion, route
changes, or variable ICMP processing.  Overall, distance establishes an
important propagation-delay floor, while routing choices and short-term
network conditions account for much of the scatter and the unusually high
maximum RTT values.

\section{Latency Breakdown}

\subsection{Method}

Five destinations were selected uniformly at random (seed 42, for
reproducibility) from the full iperf3 server list.

\begin{table}[H]
    \centering
    \caption{Five destinations selected for traceroute analysis.}
    \label{tab:traceroute-destinations}
    \begin{tabular}{lll}
        \toprule
        Destination IP & Site/Country & Outcome \\
        \midrule
        185.152.67.2 & Los Angeles, US & Non-responsive after retry (skipped) \\
        89.187.160.1 & Tokyo, JP & Reached, 13 hops \\
        105.235.237.2 & Bata, GQ & Reached, 22 hops \\
        103.146.200.98 & Port Moresby, PG & Reached, 22 hops \\
        speed2.fiberby.dk & Copenhagen, DK & Reached, 15 hops \\
        \bottomrule
    \end{tabular}
\end{table}

Traceroute was run in ICMP-echo mode (\texttt{-I}), sending 3 probes per hop
with a 2-second per-probe timeout, up to 30 hops. Per-hop RTT is the
\emph{minimum} of the responsive probes at that hop -- this approximates the
hop's propagation/processing floor and is less skewed by one transient
queuing delay than an average would be. A hop where all 3 probes time out
(shown as \texttt{* * *}) has no RTT value; it is still counted toward hop
count but excluded from the per-hop RTT data, since traceroute gives us
nothing to measure there. A destination itself is considered non-responsive,
and skipped, only if a full traceroute attempt is retried once and still
never receives a reply from the destination IP -- this distinguishes a
genuinely dead/filtered destination from a one-off transient failure (both
were observed during testing; see repo \texttt{DECISIONS.md} for detail).
Initial testing with the default UDP-mode traceroute under-reported
reachability: 3 of the first 5 randomly-selected destinations appeared
non-responsive, yet a plain \texttt{ping} showed all 3 were alive -- they (or
a router along the path) silently drop the ICMP "port unreachable" reply
UDP-mode traceroute depends on, while answering ICMP echo normally. Switching
to ICMP-mode traceroute resolved this.

The traceroute collection and parsing code is available here:
\codelink{traceroute invocation, parsing, and reachability check}{assignment1/part2/traceroute\_runner.py\#L47-L153},
with the retry/skip logic here:
\codelink{batch runner with retry-once-before-skip}{assignment1/part2/batch\_traceroute.py\#L21-L37}.
Target selection is here:
\codelink{seeded random target selection}{assignment1/part2/select\_targets.py\#L1-L76}.

\subsection{Per-Hop Latency Breakdown}

\begin{figure}[H]
    \centering
    \plotplaceholder{figures/hop_latency_breakdown.pdf}{Per-hop latency breakdown}
    \caption{Stacked latency breakdown for the responsive hops along the paths
    to the five selected destinations.}
    \label{fig:hop-breakdown}
\end{figure}

The plotting code is available here:
\codelink{stacked per-hop latency chart}{assignment1/part2/plotting.py\#L47-L77},
including how incremental (not cumulative) latency is computed:
\codelink{incremental\_latencies()}{assignment1/part2/plotting.py\#L25-L45}.

\subsection{Hop Count versus RTT}

\begin{figure}[H]
    \centering
    \plotplaceholder{figures/hop_count_vs_rtt.pdf}{Hop count versus RTT}
    \caption{Hop count versus destination RTT. Each point represents one of the
    five selected destination IP addresses.}
    \label{fig:hops-rtt}
\end{figure}

The plotting code is available here:
\codelink{hop-count-versus-RTT plot}{assignment1/part2/plotting.py\#L78-L97}.

\subsection{Discussion}

\textbf{Pattern in the stacked bar chart.} Across all four reached
destinations, the first four hops -- our local network and ISP egress in
West Lafayette -- are uniformly cheap (under 25\,ms) regardless of which
destination is being traced, since that segment of the path is identical for
every trace. Latency then accumulates almost entirely in one or two
``long-haul'' hops rather than being spread evenly across the path. For
example, the trace to Tokyo (89.187.160.1) jumps from 27\,ms at hop 9 to
190\,ms at hop 10 -- a single 163\,ms increment, roughly 84\% of the entire
194\,ms path RTT, almost certainly the trans-Pacific submarine link. The
Copenhagen trace shows the same shape at hop 10 (a jump from the last known
21\,ms up to 133\,ms), consistent with a trans-Atlantic hop. This confirms
that a stacked-bar breakdown is genuinely more informative than the raw
scatter alone: it shows \emph{where} along the path the latency is coming
from, not just the total.

Several hops in every trace returned \texttt{* * *} (fully non-responsive,
most often hops 5--8). These are very likely backbone/peering routers that
filter ICMP probes but still forward traffic normally -- the traceroute
still completes to the destination in every reached case. Because these
hops have no RTT of their own, whatever delay they actually contributed is
unavoidably folded into the next \emph{responsive} hop's incremental value
(see \texttt{incremental\_latencies()} above) -- a real limitation of
traceroute-based measurement, not an error in the analysis, and part of why
a single hop's bar segment can look larger than the physical link itself
would suggest.

\textbf{Hop count versus RTT.} The scatter plot shows a rough positive trend
-- more hops generally means more RTT, since a longer hop count is usually a
signature of a physically longer path -- but the relationship is weak, not
deterministic. Bata, GQ (105.235.237.2) and Port Moresby, PG
(103.146.200.98) both take 22 hops, yet their RTTs differ by roughly 50\,ms
(364\,ms vs. 312\,ms) -- a difference in specific routing (the Bata trace
detours through several US-based addresses before reaching Africa) and link
quality rather than hop count. Conversely, Tokyo (89.187.160.1) reaches in
only 13 hops but has a higher RTT (195\,ms) than the 15-hop Copenhagen trace
(137\,ms), because it is physically farther and crosses a longer undersea
link. In short: hop count is a weak proxy for path length, but the
\emph{physical distance and quality of the specific long-haul links
traversed} -- not the number of hops -- is what actually determines RTT.
This also mirrors what min/max RTT indicated in Part 1: RTT is set mainly by
propagation delay over distance, with hop count and instantaneous network
state (congestion, queuing at a particular router) contributing second-order
variation on top of that floor.

\end{document}
